% SPDX-License-Identifier: Apache-2.0
%
% The flagship. Built by //docs:flagship. Every listing is included
% from the tree.
\documentclass[journal,9pt]{IEEEtran}

\newcommand{\listingsize}{\fontsize{6}{7}\selectfont}
% SPDX-License-Identifier: Apache-2.0
%
% The house preamble, shared by every document in this directory. Loaded
% after \documentclass, before \title. Keeping it in one file is what
% stops two articles drifting apart in font, listing style or figure
% style.
% Computer Modern. IEEEtran selects Times (ptm) by default, so the face has
% to be chosen explicitly.
%
% Latin Modern is the Computer Modern variety used here, and the choice is
% forced rather than aesthetic. Under T1 encoding, plain `cmr` has no Type 1
% outlines unless cm-super is installed, and it is not in the pinned
% distribution, so pdflatex falls back to EC bitmaps and every glyph in the
% PDF becomes a Type 3 font. That was measured: the first build produced a
% document with 10 Type 3 fonts and no Type 1. Latin Modern is Computer
% Modern redrawn with a full T1 repertoire and real .pfb outlines, so the
% same design comes out scalable.
\usepackage[T1]{fontenc}
\usepackage[utf8]{inputenc}
\usepackage{lmodern}
% microtype is not in the pinned TeX distribution. emergencystretch alone
% keeps the two-column measure from producing overfull lines.
\setlength{\emergencystretch}{3em}

\usepackage{amsmath}
\usepackage{amssymb}
\usepackage{amsthm}
\usepackage{booktabs}
\usepackage{array}
% enumitem is not in the pinned TeX distribution; the list spacing below
% does what \setlist would have done.
\usepackage{float}
\usepackage{xcolor}
\usepackage{listings}
\usepackage{tikz}
\usetikzlibrary{arrows.meta,positioning,fit,backgrounds,calc}
\usepackage[hidelinks,bookmarks=true]{hyperref}

% Numbered, definition-style box for a rule the reader is meant to apply.
\theoremstyle{definition}
\newtheorem{rulebox}{Rule}
\newtheorem{example}{Example}

% Unnumbered asides.
\theoremstyle{remark}
\newtheorem*{tip}{Tip}
\newtheorem*{pitfall}{Pitfall}

% Tighter lists, in place of enumitem's \setlist.
\setlength{\topsep}{2pt}
\setlength{\itemsep}{1pt}
\setlength{\parsep}{0pt}

\newcommand{\code}[1]{\texttt{#1}}
\newcommand{\term}[1]{\emph{#1}}

% Syntax highlighting. Four roles, distinguishable in colour and still
% distinguishable in greyscale, because an IEEE article gets printed: the
% keyword is the only bold role, the comment is the only italic one, and the
% two remaining roles differ in lightness as well as in hue.
\definecolor{kwblue}{RGB}{20,60,150}
\definecolor{cmtgreen}{RGB}{90,120,90}
\definecolor{strred}{RGB}{150,50,40}
\definecolor{numgray}{gray}{0.45}
\definecolor{framegray}{gray}{0.70}
\definecolor{bgshade}{gray}{0.97}
% A darker fill for figure cells. The listing background has to stay very
% light to keep code readable; a figure cell has to be clearly shaded or the
% distinction it draws is invisible in print.
\definecolor{cellshade}{gray}{0.90}

% The listing size is the document's choice, because it depends on the
% column: a two-column article sets `\listingsize` to 6pt and keeps its
% listings within 70 columns, a one-column document keeps the body size
% and includes source files at 80. Lines are never broken; a line that does not fit
% overflows the frame, which is visible, rather than wrapping, which is
% not.
\providecommand{\listingsize}{\scriptsize}

% `'` is deliberately not a string delimiter: the language uses it for loop
% labels, as Rust does, so treating it as a quote would swallow the rest of
% the line.
\lstdefinestyle{txhdl}{
  basicstyle=\ttfamily\listingsize,
  keywordstyle=\bfseries\color{kwblue},
  commentstyle=\itshape\color{cmtgreen},
  stringstyle=\color{strred},
  numberstyle=\tiny\color{numgray},
  morekeywords={mod,use,pub,const,type,struct,enum,trait,impl,fn,async,let,mut,
    match,if,else,loop,while,for,in,break,continue,return,as,await,
    module,bus,channel,transaction,protocol,pipeline,flow,seq,config,
    port,stage,pipe,instance,connect,bind,tag,domain,blackbox,foreign,
    property,role,signal,handler,true,false,
    sequence,interface,modport,implements,package,import,var,wire,func,proc,
    case,default,next,fork,branch,join,state,fsm},
  morecomment=[l]{//},
  morestring=[b]",
  showstringspaces=false,
  columns=fullflexible,
  keepspaces=true,
  breaklines=false,
  % Framed, so a listing is bounded rather than floating in the column.
  frame=single,
  rulecolor=\color{framegray},
  framesep=3pt,
  backgroundcolor=\color{bgshade},
  xleftmargin=3pt,
  xrightmargin=3pt,
  aboveskip=6pt,
  belowskip=6pt,
  % Every listing is numbered and captioned, so it can be referred to by
  % number. `caption` without `float` numbers the listing and leaves it
  % where it was written, which is what a listing in running prose needs.
  captionpos=b,
  abovecaptionskip=2pt,
  belowcaptionskip=6pt,
}

% Figures drawn as text. Framed the same way, and with the highlighting
% switched off, because a box-and-arrow diagram is not code and half its
% words turning blue reads as an accident. Setting a style adds keys rather
% than clearing the ones already set, so the three roles are neutralised by
% hand rather than by leaving them out.
\lstdefinestyle{ascii}{
  basicstyle=\ttfamily\listingsize,
  keywordstyle=\ttfamily,
  commentstyle=\ttfamily,
  stringstyle=\ttfamily,
  columns=fullflexible,
  keepspaces=true,
  frame=single,
  rulecolor=\color{framegray},
  framesep=3pt,
  backgroundcolor=\color{bgshade},
  xleftmargin=3pt,
  xrightmargin=3pt,
  aboveskip=6pt,
  belowskip=6pt,
}
\lstset{style=txhdl}

% House TikZ styles. `shorten` lives on the arrow styles rather than on each
% arrow, so every arrow in the document gets the gap, including ones added
% later. TikZ clips a path at a node's border, so without it an arrowhead
% butts against the box with no gap at all. Keep `node distance` well above
% twice the shorten, or the arrow shrinks to nothing.
\tikzset{
  box/.style={draw,rounded corners=2pt,align=center,inner sep=4pt,
              font=\footnotesize},
  boxf/.style={box,fill=bgshade},
  lbl/.style={font=\scriptsize\itshape,align=center},
  fl/.style={-{Stealth[length=4pt]},shorten >=2.5pt,shorten <=2.5pt},
  fld/.style={fl,dashed},
}



% The colours of the placement map, one per subsystem, chosen to be
% told apart in print as well as on a screen. //tools/fpgamap names
% them, so they are defined here and nowhere else.
\definecolor{flagdie}{gray}{0.92}
\definecolor{flagcore}{RGB}{40,90,170}
\definecolor{flagmem}{RGB}{200,110,30}
\definecolor{flageth}{RGB}{50,130,60}
\definecolor{flagvid}{RGB}{140,60,150}
\definecolor{flagcdc}{RGB}{30,150,160}
\definecolor{flagtop}{gray}{0.45}

\InputIfFileExists{buildstamp.tex}{}{}
\providecommand{\buildstamp}{Draft build (outside the Bazel flow).}

\title{The TxHDL Flagship}

\author{Filip Filmar and Dragiša Janković%
\thanks{Both authors are independent. Filip Filmar is the
corresponding author, at f@hdlfactory.com; Dragiša Janković is
reached at linkedin.com/in/dragisa-jankovic.}%
\thanks{This document is the system the board is kept in, for a reader
who knows the core~\cite{vreteno} and the three subsystems it brings
together~\cite{eth,hdmi,axi}. Every listing is included from the
project repository \code{HDL/txhdl} at build time. The
cover~\cite{cover} lists every document.}%
\thanks{The authors used a large language model, Claude, as an
assistant in exploring the concepts and in writing and constructing
the documents and the programs; every commit of the repository says
so and carries the prompts that led to it, verbatim.}%
\thanks{\buildstamp}}

\markboth{The TxHDL Flagship}{Filmar and Janković: The TxHDL Flagship}

\begin{document}
\maketitle

\begin{abstract}
The pieces of this project are proven on the Alinx AX7A200B one at a
time: a core with a memory controller behind it, an Ethernet port that
echoes, an HDMI output that draws a picture. Each is its own bitstream,
and only one of them is in the part at a time. The flagship is the
bitstream that holds all of them, written to the board's flash, so that
the board powered on with nothing attached is the whole system rather
than whichever experiment ran last. It is also the bitstream that stops
moving: the core comes up running a loader, so the software changes in
a second over the serial port and the hardware underneath it stays
where it is. This document says what is in it, how the three clock
domains are kept apart, what the lights mean, and how it is written to
the flash.
\end{abstract}

\section{Why a flagship}
\label{sec:f-why}

A part holds one bitstream. Every demonstration in this repository is
therefore a claim about a board that is not in that state any more:
the Ethernet echo was in the part on Thursday, the HDMI picture on
Friday, and by Saturday it is the DDR3 test again. Nothing is wrong
with any of them, and the board still cannot be handed to somebody and
switched on.

The flagship answers that. One bitstream holds every part of the system
that is proven here, it lives in the board's QSPI flash rather than in
the FPGA's volatile configuration, and it is rebuilt whenever any of
its pieces changes. A board with power and nothing else is the system.

The second reason is the one that saves the most time. Place and route
of this design is half an hour. A flagship whose core ran a program
built into its netlist would cost half an hour per edit of that
program, which is why the demonstrations so far have each been their
own bitstream. The flagship's core instead comes up running the
loader~\cite{vreteno}: it waits on the serial port, takes a program,
writes it into the DDR3 memory and jumps to it. The hardware is fixed
and known; the software is whatever was sent last. Changing what the
flagship does costs a second.

\section{What is in it}
\label{sec:f-what}

Three subsystems, each unchanged from the design that proved it.

\emph{The core and its memory.} \code{board}, the netlist
\code{\#[lower]} writes for \code{vreteno32::board::Board}: the Vreteno
RV32IMC core, its AXI tracker, the router, the data memory, the timer
and software interrupt, the platform-level interrupt controller, the
pulse width modulator, the serial port, and the board's gigabyte of
DDR3 behind UberDDR3's controller. The core's boot memory holds the
loader.

\emph{The Ethernet port.} The two lowered halves of the MAC with the
PHY's RGMII around them and a crossing between the receive and
transmit clocks, echoing every good frame~\cite{eth}. The core does not
reach the port yet. The echo is what keeps the port alive and provable
from the other end of the cable while a peripheral for it is written.

\emph{The HDMI output.} The video peripheral with its framebuffer and
the I2C master that configures the SiI9134 encoder~\cite{hdmi}. Here
the flagship does add something the demonstration did not have: the
core reaches the peripheral, so a program can paint the screen.

\section{The third slot}
\label{sec:f-slot}

The core's router has five ports and all five are taken. The serial
port and the pulse width modulator already share the page at
\code{0x3000} behind an AXI-Lite bridge, each a sixteenth of it, for
exactly that reason: a peripheral of six registers does not deserve a
port of its own. The video peripheral is three registers, so it takes a
third sixteenth, at \code{0x3200}, and the bridge becomes a
\code{LiteBridge3}.

What is different about the third slot is that it does not reach a
field of \code{Board}. It leaves the unit as five channel ports:

\lstinputlisting[rangebeginprefix=//\ begin\{,rangebeginsuffix=\},%
rangeendprefix=//\ end\{,rangeendsuffix=\},includerangemarker=false,%
linerange=vslot-vslot,caption={\code{cpu/vreteno/src/board.rs}: the
third slot, and why it leaves the
unit.},label={lst:f-slot}]{cpu/vreteno/src/board.rs}

because what sits there runs on a clock of its own. The video
peripheral's raster is timed by the pixel clock, 25.2\,MHz, and the
core runs at 100\,MHz; a peripheral inside \code{Board} would have to
run on the core's clock, and a VGA raster at 100\,MHz is not a VGA
raster. So the slot is brought out, and joining it to something is the
board top's business. A board with nothing there ties it off, which is
what \code{cpu/vreteno/board/vreteno\_board.v} does; a program that
touched \code{0x3200} on such a board would wait forever, and none
does.

\section{Three clocks, one input}
\label{sec:f-clocks}

The board has one 200\,MHz differential input. The flagship buffers it
once and hangs three generators off it, each the generator its
subsystem is already proven with, so that bringing the three designs
together changes nothing about any of them.

\begin{itemize}
\item A PLL, times six to 1200\,MHz, giving 100\,MHz for the design and
  the memory controller, 400\,MHz for the memory, the same 400\,MHz a
  quarter cycle late, and 200\,MHz for the controller's delay control.
\item An MMCM, times five to 1000\,MHz over eight, giving 125\,MHz for
  the Ethernet transmit half. A second MMCM inside \code{eth\_rgmii}
  shifts the clock the PHY recovers from the line by a quarter cycle,
  because that PHY adds no receive delay of its own.
\item An MMCM, times 63 over ten and then over 50, giving 25.2\,MHz for
  the pixel clock, within half a per cent of the mode's 25.175\,MHz.
\end{itemize}

One thing crosses between the domains, and it is a channel rather than
a signal: the core's AXI-Lite to the video peripheral. Each of the five
AXI-Lite channels crosses on its own, through \code{chan\_cdc}, the
asynchronous FIFO with Gray-coded pointers that the Ethernet echo
already uses on this board. Crossing each channel on its own is the
whole of what the protocol asks, because the five have no timing
relationship to each other: a write is complete when its response
arrives, whenever that is.

\lstinputlisting[caption={\code{lib/board/chan\_cdc.v}: a lowered
unit's channel, carried from one clock to
another.},label={lst:f-cdc}]{lib/board/chan_cdc.v}

That file used to be \code{eth/board/eth\_cdc.v}, and the flagship is
what made it general: it now lives in \code{//lib/board}, which is
where Verilog that more than one area's board design wants is kept.

\section{The lights}
\label{sec:f-leds}

The board has four LEDs, lit when driven low, and the flagship has
more than four things worth watching. These four are what a person at
the board needs to tell a healthy flagship from a sick one.

\begin{itemize}
\item \textbf{LED1}, the pulse width modulator's first channel, or the
  core halted. This is the software's own light: the loader leaves it
  dark, and a program sent down the wire does with it what it likes.
\item \textbf{LED2}, a frame seen on the Ethernet port, either
  direction, held for about a fifth of a second so that a person sees
  it.
\item \textbf{LED3}, the DDR3 memory calibrated. The loader writes the
  program it receives into that memory, so a dark LED3 explains
  everything else that is wrong.
\item \textbf{LED4}, the heartbeat, three times a second. It says the
  PLL has locked and that this part holds a bitstream that is running,
  which is the first question to ask of a board that does nothing.
\end{itemize}

The HDMI output has no light of its own. The monitor is its light.

\section{The design}
\label{sec:f-top}

Everything above is one hand-written module, and it is the only new
hardware in the flagship: the pins, the three generators, the resets,
the five crossings and the lights.

\lstinputlisting[caption={\code{flagship/board/flagship.v}: the
flagship's top.},label={lst:f-top}]{flagship/board/flagship.v}

The constraints are shared rather than copied. The board's clock, its
reset, its LEDs, its serial port and its memory come from the two files
under \code{cpu/vreteno/board/}; the PHY's pins from
\code{eth/board/eth\_pins.xdc}; the encoder's from
\code{hdmi/board/hdmi\_pins.xdc}. Each pin of this board is named in
exactly one file in the tree, and every design that drives that pin
reads it, so a pin corrected for one design is corrected for the
flagship in the same change. What is left for the flagship's own file
is the part no subsystem owns: telling the three generators apart for
timing.

\lstinputlisting[caption={\code{flagship/board/ax7a200\_flagship.xdc}:
what no subsystem owns.},label={lst:f-xdc}]%
{flagship/board/ax7a200_flagship.xdc}

\section{Building it, and putting it in the flash}
\label{sec:f-build}

Every Vivado target in this repository is manual, because each is
minutes of work that a plain build should not start.
\code{bazel build //flagship:flagship\_synth} takes five minutes and
\code{//flagship:flagship\_pnr} about twenty-five more, and the second
depends on the first, so the second alone is the whole of it.

Writing the result to the board needs the tunnel to the machine the
board is attached to, which \code{bazel run
//cpu/vreteno/board/remote:hw\_server} opens. Then either
\code{//flagship:flagship\_prog} writes the FPGA itself, which the next
power cycle forgets, or \code{//flagship:flagship\_flash} writes the
QSPI flash, which is where a flagship belongs: the board then comes up
as itself with nothing attached. Both take
\code{-{}-hostport localhost:3122} and the device pattern, as the
targets' comments state.

The flash is a Micron MT25QL128 over a four-bit bus. It is the part
this board answers with, read off the chip over JTAG rather than out of
a manual, because board revisions differ here and the manual for this
one names a larger part that Vivado then refuses.

Once the flagship is in the part, a program goes to it over the serial
port and no Vivado is involved at all: \code{bazel build
//cpu/vreteno/rust:hello\_ram\_bin} for the image, and \code{bazel run
//cpu/vreteno/board/remote:load -{}-image=...} to send it. That is the
loop the flagship exists to make cheap, and it is seconds rather than
half an hour.

\section{What it costs}
\label{sec:f-cost}

The whole flagship is a small part of an xc7a200t. Synthesis reports
10\,305 slice LUTs, 7.7 per cent of the part; 7\,966 registers, 3.0 per
cent; 16 block RAM tiles of 365; four DSPs; 126 of 285 bonded pins; and
of the clocking, one PLL and three MMCMs, one of which is inside the
RGMII wrapper. The part is not the constraint, which is worth saying
plainly: the question for the next subsystem is whether it meets
timing, not whether it fits.

\section{Where it landed}
\label{sec:f-place}

Counting says how much of the part is used. It does not say where, and
where is the question a design of four subsystems raises: are they four
regions of the part or one mixture, and which of them did the pins pull
to an edge?

Figure~\ref{fig:f-map} is the answer, read off the routed checkpoint
rather than drawn by hand. \code{//flagship:place\_report} opens the
checkpoint in Vivado, attributes every placed cell to a subsystem by the
top level instance it sits under, and writes a count per subsystem per
slice column and row into \code{docs/flagship\_place.tsv}, which is
committed. \code{//tools/fpgamap} bins that and draws it at build time. The
arrangement is the one the layout numbers already use: the extraction is
manual because it is minutes of Vivado, and a Bazel action has no
repository to write into.

\begin{figure}[t]
\centering
\input{flagshipmap}
\caption{Where the flagship's subsystems landed on the xc7a200t. A
square is eight slices by eight; its colour is whichever subsystem
holds the most cells in it and its shade how full it is. The grey is
part that holds nothing.}
\label{fig:f-map}
\end{figure}

\begin{table}[t]
\caption{The same placement as numbers. A slice holds several cells,
so these counts are larger than a utilisation report's, which counts
slices.}
\label{tab:f-place}
\centering
\footnotesize
\input{flagshiptotals}
\end{table}

Three things are worth reading off it.

\emph{The subsystems are regions, not a mixture.} The Ethernet sits in
slice columns 0 to 22 and the core in 93 to 147: they do not touch. The
placer was told nothing about this. It follows from the pins, which is
the point: the PHY is wired to one bank and the memory to another, and
logic goes where its pins are.

\emph{The memory controller is the largest thing in the design}, larger
than the core, its bus and all of its peripherals together, and it is
pinned hard against the memory's bank. That is worth knowing before
adding anything that wants to be near those pins.

\emph{The crossing sits between what it crosses.} The five FIFOs that
carry the core's AXI-Lite to the pixel clock are in columns 92 to 118,
between the core and the video, which is where a crossing should be and
is not something anything arranged.

\section{The first program down the wire}
\label{sec:f-first}

The slot and the crossing are worth nothing until something writes the
screen. The first program that does is a turning icosahedron, drawn by
the core into the framebuffer with the TxHDL logo in a corner
(Listing~\ref{lst:f-ico}). It is linked for the memory at
\code{0x4000\_0000}, flattened, and sent by the loader: the picture
on the flagship changed a dozen times in an afternoon and Vivado ran
for none of them, which is the loop this document is about.

Three things in it are worth reading. The twenty faces are not a table:
they are found at startup by taking every triple of vertices that are
pairwise one edge apart and winding each so its normal points away from
the centre, and the program says how many it found, which on the board
was twenty. The face normals are unit vectors turned with the solid, so
one number decides whether a face is seen and how brightly it is lit,
and for a convex solid that test is the whole of hidden surface
removal. And nothing is cleared: a row is composed in a buffer and
written once, because the raster reads the framebuffer while the core
writes it, and a screen that is cleared and then drawn is seen in
between. The first version cleared, and it flickered.

Two things it got wrong on the way are in its comments, because they
are the kind of thing a picture on a monitor teaches and a simulation
does not: the eye was three radii from the solid, so the nearest face
swelled to nearly the whole silhouette and read as a back face painted
over the front; and each colour channel was shaded and rounded on its
own, so a face turning smoothly changed hue rather than brightness.

\lstinputlisting[caption={\code{cpu/vreteno/rust/ico\_hdmi.rs}: the
icosahedron.},label={lst:f-ico}]{cpu/vreteno/rust/ico_hdmi.rs}

The logo is data: one twelve-bit word per pixel at the size a corner
wants, in a crate of constants that builds for the core as it builds
for this machine (Listing~\ref{lst:f-logo}). It is written by
\code{//tools/logo2rs} from the PNG and committed, as the layout
numbers are committed, because a Bazel action cannot reach the drive
the logo is kept on. The dark background of the source becomes a value
that is not a colour, which a program skips, so the logo sits on
whatever is drawn under it.

\lstinputlisting[caption={\code{lib/logo/logo.rs}: the head of the
logo, as \code{logo2rs} writes it.},label={lst:f-logo},%
linerange={1-27}]{lib/logo/logo.rs}

\lstinputlisting[caption={\code{tools/logo2rs/main.go}: the PNG to the
crate.},label={lst:f-logo2rs},tabsize=2]{tools/logo2rs/main.go}

\section{What is not done}
\label{sec:f-todo}

\emph{The core does not reach the Ethernet port.} The port echoes, on
its own, beside the core. The AXI-Lite peripheral for the MAC
exists~\cite{eth}; what is missing is a sixth slot to put it on and the
crossing to the 125\,MHz domain, which is the same crossing the video
already uses.

\emph{One program paints the screen, and it is the whole of what is
known about painting it.} Section~\ref{sec:f-first} says what it
does not do: no second buffer, so a moving edge tears; no depth
buffer, which a convex solid does not need and the next one will.

\emph{There is one flagship.} The issue that asks for this one allows
for several, named \code{flagship-0}, \code{flagship-1} and so on, if
the part ever fills up. It has not: see
Section~\ref{sec:f-cost}.

\begin{thebibliography}{9}

\bibitem{vreteno}
F.~Filmar and D.~Janković, ``Vreteno in TxHDL,'' project repository
\code{HDL/txhdl}, \code{//docs:vreteno}, 2026.

\bibitem{eth}
F.~Filmar and D.~Janković, ``Ethernet in TxHDL,'' project repository
\code{HDL/txhdl}, \code{//docs:eth}, 2026.

\bibitem{hdmi}
F.~Filmar and D.~Janković, ``HDMI in TxHDL,'' project repository
\code{HDL/txhdl}, \code{//docs:hdmi}, 2026.

\bibitem{axi}
F.~Filmar and D.~Janković, ``AXI in TxHDL,'' project repository
\code{HDL/txhdl}, \code{//docs:axi}, 2026.

\bibitem{cover}
F.~Filmar and D.~Janković, ``TxHDL documents,'' project repository
\code{HDL/txhdl}, \code{//docs:cover}, 2026.

\end{thebibliography}

\end{document}
